% \begin{wrapfigure}[15]{r}{0.7\textwidth}
%   \begin{center}
%     \includegraphics[width=0.6\textwidth]{assets/antarctic_oscillation.pdf}
%   \end{center}
%   \caption{Subgraph of G depicting learned causal relationships among Antarctic Oscillation regions. The left figure illustrates the divisions between merged nodes, while the right figure visualizes the spatial factors and causal links.}
%   \label{fig:antarctic}
% \end{wrapfigure}

% \begin{figure}[t]
%     \centering
%     \includegraphics[width=0.85\columnwidth]{assets/temp_subgraph.jpg}
%     \caption{Qualitative results for Global Temperature climate dataset. The numbers on the arrow refers to the time lag of the causal links. Subgraph of G depicting learned causal relationships among regions associated with the (a) Northern Atlantic Oscillation (b) Antarctic Oscillation.}
%     \label{fig:temp_subgraph}
% \end{figure}


\begin{figure*}[t]
    \centering
    \includegraphics[width=0.7\textwidth]{assets/G_subgraphs_MJO.pdf}
    \includegraphics[width=0.7\textwidth]{assets/G_subgraphs_NAO.pdf}
    \includegraphics[width=0.7\textwidth]{assets/G_subgraphs_AAO.pdf}
    \caption{\textbf{Visualization of discovered causal relationships in climate datasets.} Subgraph among regions associated with the (top) Madden-Julian Oscillation, (middle) Northern Atlantic Oscillation, (bottom) Antarctic Oscillation. Causal links are shown for (left) Precipitation, (middle) Temperature, and (right) cross-variate interactions. Numbers on the edges indicate time-lag, while the numbers on the nodes indicate the node indices. }
    \vspace{-3mm}
    \label{fig:all_osc_subgraphs}
\end{figure*}